% LNCS: tóm tắt 150--250 từ (tính trên bản tiếng Anh). Không đưa giá trị p vào tóm tắt.
% Số: facts.md mục Test-4 (40, 6, 2 họ; cặp lệch 4/0; cùng chiều trên ba mô hình; ΔU −3,1; token và độ trễ v3 so với v1),
% thích nghi, AgentDojo v3. 20/11 (góp ý owner): nêu kết quả so với luật so khớp chuỗi theo hướng tích cực nhưng trung
% thực, H1 không đạt được diễn đạt là "cần tập kiểm thử lớn hơn để khẳng định bằng kiểm định đăng ký trước".
% 20/11 (soát văn phong): tách câu dài, bớt "còn", "cưỡng chế" thành "thực thi chính sách".
\begin{abstract}
Một chỉ thị cài vào nội dung mà tác tử LLM đọc qua công cụ có thể khiến tác tử gửi dữ liệu tới đích lạ, làm lộ giá trị
mật hoặc sửa tham số của một hành động hợp lệ. Bài báo đề xuất \ours, một cơ chế thực thi chính sách tất định đặt tại Tool
Broker, gồm chính sách ý định có xét phủ định trích từ yêu cầu của người dùng, nhãn độ nhạy cảm cho từng giá trị và phép
ghim các giá trị người dùng đã nêu, không gọi thêm LLM. Ba thuộc tính cấp quyền, bảo mật và toàn vẹn được chứng minh trên
một mô hình chuyển trạng thái dưới năm giả định nêu rõ. Trên tập kiểm thử tiếng Việt gồm 48 họ tấn công, niêm phong trước
khi viết mã, với ba mô hình nền mở trọng số, \ours{} giảm số họ bị tấn công thành công từ 40 khi không phòng thủ xuống 2.
Khi ba thành phần được thay bằng luật so khớp chuỗi, còn 6 họ bị hại. \ours{} chặn thêm bốn họ trong số đó, mỗi họ bởi
đúng thành phần được thiết kế cho nó, không để lọt họ mới, và cho cùng chiều hiệu ứng trên cả ba mô hình. Do luật so khớp
chuỗi đã gần mức sàn, cần một tập kiểm thử lớn hơn để khẳng định khác biệt này bằng kiểm định đăng ký trước. Tiện ích chỉ
giảm 3,1 điểm phần trăm so với khi không phòng thủ, còn token và độ trễ chênh không quá 1,4\% so với luật so khớp chuỗi.
Dưới hai vòng tấn công thích nghi, \ours{} không để lọt họ nào ngoài những họ đã lọt với tấn công gốc. Trên AgentDojo
bằng tiếng Anh và tiếng Việt, số tấn công thành công được giữ ở mức thấp, đổi lại tiện ích giảm ở các nhiệm vụ lấy đích
từ dữ liệu công cụ.

\keywords{Chèn lệnh gián tiếp \and Tác tử LLM \and Cấp quyền theo ý định \and Kiểm soát luồng thông tin \and Tiếng
Việt}
\end{abstract}
